% Supplementary Notes 1 to 5 (Task 46): v3 Appendices A to D and the limitation detail that the Discussion of the main
% text does not keep (plan 42, section 1.5). Every result number is a macro. Supplementary items are referenced with
% \ref inside this document and main-text items with the numbering macros of generated/srep_numbers.tex.

\section*{Supplementary Note \suppNoteImpl: Implementation details}
\label{sec:si_impl}
Supplementary Table~\ref{t_impl} lists the settings used on both datasets.
The FT-Transformer was trained with AdamW, a linear warm-up followed by cosine decay of the learning rate, gradient-norm
clipping, class weights proportional to the square root of the inverse class frequency (\apClassWeight), bf16
autocast during training and float32 for evaluation and attribution, and early stopping on validation macro-F1.
On 2018 all three seeds ran to the maximum of \apMaxEpochs{} epochs without early stopping (best epochs
\apBestEpochEighteen), and the validation macro-F1 ranged from \valFlatEighteenRange{} over the last
\apPatience{} epochs of every seed; on 2017 early stopping ended the runs after \apEpochsRunSeventeen{} epochs.
As a diagnostic of the attention explanations we selected, per seed, the single attention head (the classification-token
row of one of the \apHeads{} heads in one of the \apBlocks{} layers) with the lowest deletion AUC on the validation
explanation rows, macro-averaged over attack classes and with features masked by zero in the standardised space, and
evaluated it on the test explanation rows with training-mean masking like the other methods.
Features were log-transformed when non-negative with training skewness above the threshold, standardised with training
statistics and clipped.
Attacks used sign-gradient projected gradient descent with a cosine-decayed step size given as a fraction of the
per-feature budget, the best projected iterate per row over all steps and restarts, and one random restart.
The \nConsistency{} statistics of control c2 are the softmax entropy, energy, logit norm, the two largest logits and the
gap between the second and third, the benign logit, the largest attack-class logit, the logit standard deviation, the
benign probability, and the Mahalanobis and \paramKnn{}-nearest-neighbour distances for the predicted class.
Detector D used a very large inverse regularisation strength (effectively unregularised); D$_+$ used the value in the
table.
For a1 and a2 the draw was repeated with a second shuffle seed, and regenerating the training rows, the scores and the
A6 rows with the original random streams reproduced D, D$_+$ and the A6 rows exactly.
A8x computes IG in the loss with the same Gauss-Legendre rule as the true scores and was checked against the captum
implementation and against the saved true scores before any attack (Supplementary Table~\ref{t_runs}); A6x runs over
all test source rows in chunks of the size used by A6, so that its random restart and batch shapes equal those of A6.
The datasets, their classes and splits are in Supplementary Table~\ref{t_data}, and the constraint model in
Supplementary Table~\ref{t_constraints}.

\section*{Supplementary Note \suppNoteRand: Parameter randomisation check}
\label{sec:si_rand}
As a sanity check in the sense of Adebayo et al. \cite{adebayo2018}, we re-initialised parts of each trained
FT-Transformer at \nRandLevels{} levels (the classifier head alone; the transformer blocks cascading down from the last
block, and then also the feature tokenizer; and each block alone, with the classifier head kept) and recomputed the
three attention variants and IG on the first \randRowsPerClass{} rows per class of the clean explanation set.
Supplementary Table~\ref{t_randomisation} gives, for three of these levels, the Spearman correlation and the
top-\paramTopK{} Jaccard overlap between the explanation of the original and of the randomised model.
With all blocks and the tokenizer re-initialised, the rollout ranking was unrelated to the original one (Spearman
\randRolloutAllSeventeen{} on 2017 and \randRolloutAllEighteen{} on 2018), and so was the signed IG
(\randIGSignedAllSeventeen{} and \randIGSignedAllEighteen), while the absolute IG kept a Spearman of
\randIGAbsAllSeventeen{} and \randIGAbsAllEighteen.
Re-initialising only the last block left rollout almost unchanged (\randRolloutTopOneSeventeen{} and
\randRolloutTopOneEighteen), which is expected because rollout also combines the attention of the unchanged lower
blocks, whereas the last-layer attention was no longer related to the original (\randLastTopOneSeventeen{} and
\randLastTopOneEighteen).
The attention explanations and the signed IG therefore changed when the learned weights were replaced; no
label-randomisation check was run.

\section*{Supplementary Note \suppNoteNeardup: Near-duplicate analysis}
\label{sec:si_neardup}
For every test row we computed the exact Euclidean distance to its nearest training row in the standardised feature
space of its dataset and flagged rows with a distance below \paramNDThr.
These flags covered \ndSlowlorisSeventeen{} of the DoS Slowloris and \ndWebSeventeen{} of the web attack test rows on
2017, \ndSlowlorisEighteen{} of the DoS Slowloris rows on 2018, and \ndAllSeventeen{} and \ndAllEighteen{} of all test
rows.
The 2018 attack source rows excluded flagged DoS Slowloris rows by design, while the 2017 attack and detector source
rows were drawn without this filter: on the \nCommonClasses{} common classes they contained \ndSrcCommonSeventeen{}
flagged rows per seed.
Because the 2018 source rows had already been filtered for near duplicates, we decided before computing V1 that
cross-dataset comparisons should use rows without near duplicates; we fixed V1 as the primary set only after seeing the
sensitivity results.
Supplementary Table~\ref{t_limitations_data} lists the small classes and the share of test rows with a near duplicate
in the training set, and Supplementary Table~\ref{t_limitations_stages} the flagged source rows and benign test
negatives of the attack and detector experiments and the seeds of every pipeline step.
Supplementary Table~\ref{t_neardup_split} compares rows with and without a near duplicate.
On 2017, flagged rows were easier for both detectors: on A6 rows ($\lambda = 1$) the TPR of D was \tprDNdSeventeen{} on
flagged and \tprDOtherSeventeen{} on other source rows (D$_+$: \tprDplusNdSeventeen{} and \tprDplusOtherSeventeen), and
the FPR of D was \fprDNdSeventeen{} on flagged and \fprDOtherSeventeen{} on other benign test negatives (D$_+$:
\fprDplusNdSeventeen{} and \fprDplusOtherSeventeen).
On 2018 no attack source row was flagged, and the FPR split was mixed: at $\epsilon = 1$, D \fprDNdEighteen{} on flagged
and \fprDOtherEighteen{} on other negatives (D$_+$: \fprDplusNdEighteen{} and \fprDplusOtherEighteen); at
$\epsilon = 0.5$, D \fprDNdEighteenHalf{} and \fprDOtherEighteenHalf{} (D$_+$: \fprDplusNdEighteenHalf{} and
\fprDplusOtherEighteenHalf).
Flagged rows made up \ndNegSeventeen{} and \ndNegEighteen{} of the benign test negatives.
Supplementary Tables~\ref{t_sensitivity_2017_attacks} to~\ref{t_sensitivity_2018_defenders} give every attack and
detector quantity on V0, V1 and V2; their detector rows use D$_+$ against A8, the stand-in attacker.
Moving from V0 to V1 and V2 kept the levels of successful and undetected rows in their order of magnitude on both
datasets (2017 D under A6: \undetDAsixOneVzeroSeventeen, \undetDAsixOneVoneSeventeen, \undetDAsixOneVtwoSeventeen).

\section*{Supplementary Note \suppNoteResults: Supplementary results}
\label{sec:si_results}
\paragraph{Attacks by budget and all source rows.}
Supplementary Fig.~\ref{f_evasion_eps} and Supplementary Table~\ref{t_evasion} give constrained evasion success by
budget, and Supplementary Tables~\ref{t_clean_agreement}, \ref{t_clean_faith} and~\ref{t_faith_attack} the clean-data
and attack tables of the main text on all rows (V0) next to V1, with the A3 variants and budgets that the main text
omits.
Supplementary Fig.~\ref{f_mass_ratio} shows the attention mass ratio of Table~\mainTabAttn{} of the main text, panel a.
Supplementary Tables~\ref{t_defenders} and~\ref{t_per_class} give the detector results per attack and per class on all
source rows (V0), and Supplementary Figs.~\ref{f_detection_tpr} and~\ref{f_per_class_heatmap} the corresponding TPR and
per-class union; in these, D$_+$ faces the stand-in attacker A8.
Supplementary Table~\ref{t_per_class_v1} gives the per-class shares on V1 behind the class statements of the main text.

\paragraph{Clean-data results that did not replicate.}
Three results on 2017 did not hold in the same direction on 2018 (all explanation-set rows): the advantage of a static
class-mean rollout ranking over the per-row ranking (AOPC gap of static minus per-row rollout \staticGapSeventeen{},
seed range \staticGapSeventeenRange, against \staticGapEighteen, \staticGapEighteenRange); the advantage of $|x|$ over
rollout (deletion rank \rankDelAbsXSeventeen{} and \rankDelAbsXEighteen{} for $|x|$, \rankDelRolloutSeventeen{} and
\rankDelRolloutEighteen{} for rollout); and the advantage of the single attention head selected on validation data
(difference to rollout \headMinusRolloutSeventeen, seed range \headMinusRolloutSeventeenRange, against
\headMinusRolloutEighteen, \headMinusRolloutEighteenRange).
The method orderings on V0 and V2 correlated across the datasets with \ordSpDelVzero{} and \ordSpDelVtwo{} (deletion)
and \ordSpInsVzero{} and \ordSpInsVtwo{} (insertion).

\paragraph{Against the stand-in attacker A8.}
Supplementary Table~\ref{t_union_a8} gives the equal-effort union with D$_+$ facing A8 on all row sets, and the
comparison of each control with D$_+$ under A8.
The 2017 difference D$_+$ minus D under A8 was \axStandinVzeroSeventeen{} (\axStandinVzeroSeventeenCI) on V0,
\axStandinVoneSeventeen{} (\axStandinVoneSeventeenCI) on V1 and \axStandinVtwoSeventeen{} (\axStandinVtwoSeventeenCI)
on V2, and the paired TPR difference on A8 rows changed sign with the row set for some settings (for $\lambda = 100$:
\tprAeightHundredVzeroSeventeen, \tprAeightHundredVoneSeventeen, \tprAeightHundredVtwoSeventeen).
The control columns of Supplementary Table~\ref{t_union_a8} compare each control with D$_+$ against the stand-in
attacker A8; they are not a comparison at the stronger attacker, which is in Table~\mainTabUnion{} of the main text.
The stand-ins reproduced the true consistency scores with a median $R^2$ of \standinMedSeventeenRange{} on 2017 and
\standinMedEighteenRange{} on 2018 across seeds; Supplementary Table~\ref{t_standin} gives every score for A8 and for
the attackers of a1 and a2.

\paragraph{A8x and A6x.}
Supplementary Table~\ref{t_a8x_setting} gives D$_+$ minus D under A8x per attack setting, the unequal-effort
sensitivity unions and the A6x union, and Supplementary Table~\ref{t_a8x_tpr} the paired TPR and AUROC differences on
the A8x and A8 rows on every row set.
With the A8 and A8x candidates together against D$_+$, an unequal-effort sensitivity union that favours the attacker of
D$_+$, the 2017 difference was \axSensVoneSeventeen{} (\axSensVoneSeventeenCI); with A6 and A6x together against D it
was \sxSensVoneSeventeen{} (\sxSensVoneSeventeenCI).
D's undetected share under A1 and A6x equalled that under A1 and A6 in every seed and row set (largest absolute
difference \sxSelMaxAbs).
Supplementary Table~\ref{t_runs} records the runs: rows, invalid rows, time, memory and the sanity checks that compared
the quantities computed inside A8x and A6x with the saved values before any attack.
The first sanity run of A8x failed at tolerance $10^{-5}$ on the Mahalanobis score and the D$_+$ decision, from
single-precision noise of the batch shape; the tolerance was set to $10^{-3}$ for these two quantities and the
\axMargin{} selection margin added before any pilot or attack, and the margin changed the in-loop count of
undetected rows from \axInloopNoMarginSeventeen{} to \axInloopMarginSeventeen{} on 2017 and no evaluated row.
The first sanity run of A6x, at a chunk smaller than that of A6, changed the path of \sxRunoneRows{} of the checked rows;
at the chunk of A6 it reproduced the saved A6 rows exactly, and A6x was run at that chunk.

\section*{Supplementary Note \suppNoteLimits: Further limitations}
\label{sec:si_limits}
\paragraph{Constraint model.}
The split into \nControls{} controls, \nDerived{} derived and \nImmutable{} immutable features, the budget rule and the
class-family freezes were designed by hand from feature definitions, and other reasonable choices could change attack
success and the attention measurements.
At $\epsilon = \paramEpsHead$ a control may change by up to its full clean value, and the random control A0 already
flipped \azeroBenignShareVzeroSeventeen{} (2017, V0) and \azeroBenignShareVzeroEighteen{} (2018) of the row
observations to benign, so the results at $\epsilon = 0.5$ and $\epsilon = \paramEpsHead$ should be read together.
We did not attack IG at an equal budget, so the stability of attention is not compared with that of an attribution
method.
The attackers knew the model weights and, for A6, A6x, A8, A8x and the attackers of the controls, the detector; this is
a strong assumption for evasion and a reasonable one for evaluating detectors, but results for weaker attackers may
differ.
The attackers of all controls kept the selection rule of A6; A6x showed that this rule and the rule of A8x give the
same flipped and undetected status for D, which was not tested for the controls.
The attackers of a1 and a2 used stand-ins with a median $R^2$ of \rtwoMedAoneSeventeen{} and \rtwoMedAtwoSeventeen{} on
2017 and \rtwoMedAoneEighteen{} and \rtwoMedAtwoEighteen{} on 2018.

\paragraph{Detectors.}
D is a linear model on four scores and D$_+$ adds \nConsistency{} linear terms; non-linear detectors might use the same
scores differently.
Control a1 uses the true class of every row, also at test time, and keeps the class-conditional distribution of
attention; its realised test FPR on 2017 (\ablFprAoneSeventeen) was higher than that of D$_+$
(\ablFprDplusSeventeen; V1, $\epsilon = \paramEpsHead$).
The validation AUROC of every detector is in-sample (D \valAurocDSeventeenRange{} on 2017 and \valAurocDEighteenRange{}
on 2018 over seeds; D$_+$ \valAurocDplusSeventeenRange{} and \valAurocDplusEighteenRange), and the realised test FPR of
D was \testFprDSeventeenRange{} and \testFprDEighteenRange.
The consistency scores of D$_+$ include IG, which costs more than attention (on 2017, \msPerRowIGRange{} ms per row
for IG with \apIGSteps{} steps against \msPerRowAttnRange{} ms for attention); we did not measure the end-to-end
latency of the detectors.

\paragraph{Data.}
The 2018 data have no timestamps, so we split at random.
Test rows within \paramNDThr{} of a training row were common in both datasets; the filter was applied after the
experiments, and models and detectors were trained with these rows present.
Web attack had \nWebTestEighteen{} test rows and \nWebValEighteen{} validation explanation rows on 2018, so its
per-class estimates are uncertain; FTP-Patator and DoS Slowhttptest were removed from 2018 and PortScan does not occur
in it, so all comparisons use the \nCommonClasses{} common classes.
The 2018 models ran to the maximum number of epochs without early stopping, so they may not have converged
(Supplementary Note~\suppNoteImpl).
Known design and error issues exist in benchmark NIDS datasets more generally \cite{flood2024,lanvin2022}.
Faithfulness was measured by deletion and insertion curves, not by human judgement.
